\section{Buchberger's criterion}\label{sec:buchberger-criterion}

\begin{theorem}[Buchberger's criterion]
\label{thm:buchberger-criterion}
Assume Condition~\eqref{eq:unit-condition} and assume that $R$ has no zero
divisors.  Then
\[
  \ideal{\LT(f)\mid f\in\ideal{G}}
  =\ideal{\LT(g)\mid g\in G}
  \quad\Longleftrightarrow\quad
  \forall p,q\in G,\;\spoly(p,q)\redstar{G}0.
\]
\end{theorem}

The quantification in Theorem~\ref{thm:buchberger-criterion} includes zero and
diagonal pairs.  These cases add no substantive condition because
\[
  \spoly(0,q)=\spoly(p,0)=\spoly(p,p)=0,
\]
and the reflexive closure gives \(0\redstar{G}0\).  Moreover, every
polynomial that actually performs a one-step reduction is nonzero by
Definition~\ref{def:one-step-reduction}.

For the forward implication, every S-polynomial of two elements of $G$
belongs to $\ideal{G}$, so Theorem~\ref{thm:groebner-characterizations}
shows that it reduces to zero.  For the converse implication, assume that all
S-polynomials of pairs in $G$ reduce to zero.  The proof establishes local
confluence by analyzing a peak
\[
  f\red{G}f_1,
  \qquad
  f\red{G}f_2.
\]
Write the two reductions as $f_i=f-q_i$, where
$q_i=c_i x^{s_i} p_i$ cancels the term with exponent $t_i$.  If the two steps
cancel the same exponent of $f$, then there exist an exponent $v$ and a
coefficient $k\in R$ such that
\[
  f_2-f_1=q_1-q_2=kx^v\spoly(p_1,p_2).
\]
Finite reduction is preserved under multiplication by a monomial term, so the
hypothesis $\spoly(p_1,p_2)\redstar{G}0$ implies that
$f_2-f_1\redstar{G}0$.

If the two steps cancel different exponents, order the cancelled monomials,
say $x^{t_1}<_m x^{t_2}$.  The polynomial $q_1$ has leading exponent $t_1$,
so its coefficient at $t_2$ is zero.  Hence $q_1-q_2$ reduces in one step to
$q_1$, and $q_1=c_1x^{s_1} p_1$ reduces to zero.  The case
$x^{t_2}<_m x^{t_1}$ is symmetric.  Since
$f_2-f_1=q_1-q_2\redstar{G}0$, the difference-to-joinability part of the
polynomial translation lemma yields $f_1\joinred{G}f_2$.  Hence every one-step
peak is joinable, so $\red{G}$ is locally confluent.  The local-confluence equivalence from
Section~\ref{sec:groebner-criterion} then proves
\lean{MonomialOrder.IsGroebner}.

Listing~\ref{lst:buchberger-criterion} states both nontrivial assumptions:
every nonzero $g\in G$ has unit leading coefficient, and $R$ has no zero
divisors.  The no-zero-divisor assumption is used when computing leading terms
of polynomial products.

\begin{lstlisting}[style=lean,caption={Buchberger's criterion for the self-relative Gr\"obner predicate.},label={lst:buchberger-criterion}]
theorem IsGroebner.iff_sPolynomial_reflTransGen_zero
    [NoZeroDivisors R]
    (G : Set (MvPolynomial σ R))
    (hG₀ :
      ∀ g ∈ G,
        IsUnit (m.leadingCoeff g) ∨ g = 0) :
    m.IsGroebner G ↔
      ∀ g₁ ∈ G, ∀ g₂ ∈ G,
        m.sPolynomial g₁ g₂ ⟶*[m, G] 0
\end{lstlisting}

The file also contains an ideal-relative form.  If
$\operatorname{Ideal.span}(G)=I$, then
\lean{MonomialOrder.IsGroebnerBasis} is equivalent to the statement that every
S-polynomial of two elements of $G$ reduces to zero modulo $G$.
